\section*{الفصل \ref{ch:b2:crustal-melting} --- انصهار القشرة: الميغماتيت والغرانيت}

\begin{solution}{exo:b2:crustal-melting:1}
(أ) لكل \qty{100}{\gram}: $n(\ce{Al2O3}) = 17.03/101.96 = \qty{0.16703}{mol}$،
$n(\ce{CaO}) = 0.09183$، $n(\ce{Na2O}) = 0.06783$، $n(\ce{K2O}) = 0.03076$.
\mbox{\textdir TLT A/CNK $= 0.16703/0.19042 = 0.877$}؛ \mbox{\textdir TLT A/NK $= 0.16703/0.09859 = 1.694$}.
(ب) \mbox{\textdir TLT A/CNK $< 1 <$ A/NK}: \omterm{def:b2:crustal-melting:asi}{ميتاألوميني}. والكالسيوم الذي لا تستطيع الفلسبارات أخذه يذهب إلى
سيليكات الكالسيوم: الهورنبلند (أو الكلينوبيروكسين، أو التيتانيت).
\end{solution}

\begin{solution}{exo:b2:crustal-melting:2}
(أ) يحمل المسكوفيت والغارنت ألومنيومًا زائدًا عن الفلسبارات: \omterm{def:b2:crustal-melting:asi}{فوق ألوميني}؛ ويمكن أن يكون هذا \omterm{def:b2:crustal-melting:types}{الغرانيت
من النمط S}. (ب) يحمل الهورنبلند والتيتانيت كالسيومًا زائدًا عن الفلسبارات:
\omterm{def:b2:crustal-melting:asi}{ميتاألوميني}. (ج) يحمل الإيجيرين والأمفيبول الصودي صوديومًا دون ألومنيوم
الفلسبار: \omterm{def:b2:crustal-melting:asi}{فوق قلوي}.
\end{solution}

\begin{solution}{exo:b2:crustal-melting:3}
كيلومتر واحد من الصخر يزن $2800\times9.81\times1000 = \qty{2.747e7}{Pa}$، ومنه
$P = 0.2747\,z$ (بوحدة \unit{kbar}، و$z$ بوحدة \unit{km}). وخط الصلابة هو $T = 680 - 7.5(P -
2) = 695 - 2.060\,z$. وبوضع $10 + 20z = 695 - 2.060z$ نحصل على $z = 685/22.060 =
\qty{31.1}{km}$، $T = \qty{631}{\celsius}$، $P = \qty{8.5}{kbar}$.
\end{solution}

\begin{solution}{exo:b2:crustal-melting:4}
(أ) الطبقات البيضاء ليوكوسومات، والأشرطة السوداء ميلانوسومات، والنيس العادي هو
الميزوسوم. (ب) ليوكوسومات وميلانوسومات وميزوسومات منفصلة ومحدّدة جيدًا: ميتاتكسيت،
ذو بنية ستروماتية. (ج) صار الصخر دياتكسيتًا (شليريًا وسديميًا):
فقد تقدّم الانصهار، وانصهرت المعادن الداكنة أيضًا، وبقيت \omterm{def:b2:magma-emplacement:magma}{الصهارة} في مكانها،
فمحت التطبّق.
\end{solution}

\begin{solution}{exo:b2:crustal-melting:5}
(أ) $M = 39.10 + 3\times24.31 + 26.98 + 3\times28.09 + 12\times16.00 + 2\times1.008 =
\qty{417.30}{\gram\per\mole}$؛ $w = 18.02/417.30 = \qty{4.32}{\percent}$.
(ب) $F = 0.25\times0.0432/0.040 = 0.270$، أي \qty{27}{\percent} من \omterm{def:b2:magma-emplacement:magma}{الصهارة} كتليًا: ضمن
20 إلى \qty{30}{\percent} حجميًا الناتجة في تجارب \omterm{def:b2:crustal-melting:reactions}{انصهار نزع الماء} من البيوتيت على
صخور مماثلة.
\end{solution}

\begin{solution}{exo:b2:crustal-melting:6}
(أ) الفلسبار البوتاسي \qty{278.35}{\gram\per\mole}، والألبيت 262.24، والأنورثيت 278.22، والمسكوفيت
398.33. (ب) الكميات: الفلسبار البوتاسي $30/278.35 = \qty{0.10778}{mol}$، والألبيت $0.13347$،
والأنورثيت $0.01797$، والمسكوفيت $0.01255$. الفلسبار البوتاسي هو \babelsublr{$\tfrac12$\ce{K2O.}$\tfrac12$\ce{Al2O3}}،
والألبيت \babelsublr{$\tfrac12$\ce{Na2O.}$\tfrac12$\ce{Al2O3}}، والأنورثيت \ce{CaO.Al2O3}، والمسكوفيت
\babelsublr{$\tfrac12$\ce{K2O.}$\tfrac32$\ce{Al2O3}} لكل وحدة صيغة. $n(\ce{Al2O3}) = 0.05389 +
0.06673 + 0.01797 + 0.01883 = \qty{0.15742}{mol}$؛ $n(\ce{K2O}) = 0.05389 + 0.00628 =
0.06017$، $n(\ce{Na2O}) = 0.06673$، $n(\ce{CaO}) = 0.01797$، والمجموع $0.14487$.
\mbox{\textdir TLT A/CNK $= 0.15742/0.14487 = 1.087$}. (ج) دون المسكوفيت تعطي الفلسبارات الثلاثة
1 بالضبط (\cref{prop:b2:crustal-melting:feldspar})؛ والمسكوفيت، الذي يحوي \ce{Al2O3} واحدًا
زائدًا لكل وحدة صيغة مقارنةً بالفلسبار البوتاسي، يجعل الصخر \omterm{def:b2:crustal-melting:asi}{فوق ألوميني}.
\end{solution}

\begin{solution}{exo:b2:crustal-melting:7}
البازلت: $900 + 2.5\times(1200 - 850) = \qty{1775}{\mega\joule\per\cubic\metre}$. والقشرة:
$2.5\times(850 - 650) + 0.20\times900 = \qty{680}{\mega\joule\per\cubic\metre}$.
$V_c/V_b = 1775/680 = 2.6$. ويتجاهل التقدير الحرارة التي تتسرّب بالتوصيل
إلى الأعلى وإلى الجوانب والزمن الذي تحتاج إليه الحرارة لتنتشر؛ ويفترض أيضًا أن كل
حرارة البازلت تبلغ قشرة مهيّأة للانصهار: فهو حدّ أعلى.
\end{solution}

\begin{solution}{exo:b2:crustal-melting:8}
(أ) $n(\ce{Al2O3}) = 14.15/101.96 = 0.13878$، $n(\ce{CaO}) = 0.01753$، $n(\ce{Na2O}) =
0.07809$، $n(\ce{K2O}) = 0.04416$؛ \mbox{\textdir TLT A/CNK $= 0.13878/0.13978 = 0.993$}، \mbox{\textdir TLT A/NK $=
0.13878/0.12225 = 1.135$}: \omterm{def:b2:crustal-melting:asi}{ميتاألوميني}، بالكاد. (ب) مع \qty{14.30}{\percent} من \ce{Al2O3}،
$0.14025/0.13978 = 1.003$؛ ومع \qty{14.00}{\percent}، $0.13731/0.13978 = 0.982$. فالمؤشر
يقع على جانبي 1 ضمن عدم اليقين المصدَّق: فالأوبسيديان يقع على
الحدّ بين \omterm{def:b2:crustal-melting:asi}{الميتاألوميني} \omterm{def:b2:crustal-melting:asi}{وفوق الألوميني}، وتصنيفه غير متين.
\end{solution}

\begin{solution}{exo:b2:crustal-melting:9}
(أ) المتوسطان $\bar P = \qty{5.0}{kbar}$، $\bar w = \qty{10.045}{\percent}$؛
$\sum(P - \bar P)(w - \bar w) = 17.62$، $\sum(P - \bar P)^2 = 14.0$؛ والميل
$17.62/14.0 = \qty{1.259}{\percent}$ لكل كيلوبار، والتقاطع $10.045 - 1.259\times5.0 =
\qty{3.75}{\percent}$. ويتوافق الميل مع القيمة المذكورة 1.27 في حدود
\qty{1}{\percent}. (ب) $3.75 + 1.259\times6 = \qty{11.30}{\percent}$. (ج) $10/0.1130 =
\qty{88.5}{\gram}$ من \omterm{def:b2:magma-emplacement:magma}{الصهارة}. (د) يستند الخط إلى بيانات بين 3 و\qty{8}{kbar} عند
\qty{800}{\celsius} فقط؛ ولا شيء يضمن أن تبقى الذوبانية خطّية عند ضغط
أعلى، حيث تتشابه \omterm{def:b2:magma-emplacement:magma}{الصهارة} والمائع أكثر فأكثر.
\end{solution}

\begin{solution}{exo:b2:crustal-melting:10}
(أ) يصعد الخط \qty{4}{kbar} على مدى \qty{50}{\celsius}: $P = 6 + 0.08\,(T - 750)$؛
عند \qty{770}{\celsius}، $P = \qty{7.6}{kbar}$. (ب) من 7.6 إلى \qty{6}{kbar} تنخفض
درجة الحرارة على امتداد الخط من 770 إلى \qty{750}{\celsius}؛ والحرارة المتحرّرة،
$2.5\times20 = \qty{50}{\mega\joule\per\cubic\metre}$، تصهر $50/900 = 0.056$، أي نحو
\qty{6}{\percent} من الصخر حجميًا. (ج) هذه هي النافذة المقترحة
لليوكوغرانيت الهيمالايي، 6 إلى \qty{8}{kbar} و750 إلى \qty{770}{\celsius}: فانخفاض الضغط
وحده يصنع صهارة، لكن بضعة في المئة فقط في كل عبور للمنحنى.
\end{solution}

\begin{solution}{exo:b2:crustal-melting:11}
$\ell^2/\kappa = (2000)^2/10^{-6} = \qty{4e12}{s} = \qty{0.13}{Myr}$ للعتبة؛
و$(10^4)^2/10^{-6} = \qty{1e14}{s} = \qty{3.2}{Myr}$ للطبقة المبطَّنة. وكلاهما قصير
بجانب عشرة ملايين سنة: فحدث واحد لا يستطيع إبقاء القشرة منصهرة كل هذه المدة، التي تحتاج إلى
حرارة تُمدّ مرة بعد مرة، بتداخلات متكرّرة أو بالتسخين البطيء لقشرة
مثخنة.
\end{solution}

\begin{solution}{exo:b2:crustal-melting:12}
(أ) $5.0x + 2.0(1 - x) = 2.75$ تعطي $x = 0.75/3.0 = 0.25$: ربع النيوسوم
ليوكوسوم. (ب) $0.25\times0.20 + 0.75\times4.0 = 0.05 + 3.00 = \qty{3.05}{\percent}$، كما
في الميزوسوم: فتفسير النظام المغلق صالح. (ج) لكل \qty{100}{\gram} من
الصخر الأم، تكوّن \qty{25}{\gram} من \omterm{def:b2:magma-emplacement:magma}{الصهارة} و\qty{75}{\gram} من البقية؛ فإذا غادر $m$ غرامًا من
\omterm{def:b2:magma-emplacement:magma}{الصهارة}، $(25 - m)/(100 - m) = 0.10$، ومنه $25 - m = 10 - 0.1m$ و$m = 15/0.9 =
\qty{16.7}{\gram}$: غادر ثلثا \omterm{def:b2:magma-emplacement:magma}{الصهارة} المكشف.
\end{solution}

\begin{solution}{pb:b2:crustal-melting:1}
\noindent\textbf{الجزء الأول.}

\textbf{1.} $39.10 + 3\times26.98 + 3\times28.09 + 12\times16.00 + 2\times1.008 =
\qty{398.33}{\gram\per\mole}$.

\textbf{2.} $18.02/398.33 = \qty{4.52}{\percent}$.

\textbf{3.} $200\times0.0452 = \qty{9.04}{\gram}$ من الماء لكل كيلوغرام من الشيست.

\textbf{4.} $9.97 + (13.82 - 9.97)\times2/3 = \qty{12.54}{\percent}$. وأقل صهارة هي
$9.04/0.1254 = \qty{72.1}{\gram}$ لكل كيلوغرام: $F \ge \qty{7.21}{\percent}$
(\cref{prop:b2:crustal-melting:water}).

\textbf{5.} $9.04/0.060 = \qty{151}{\gram}$: $F = \qty{15.1}{\percent}$.

\textbf{6.} مثلًا: قد يحوي الصخر مسكوفيتًا أقل؛ وقد يكون جزء من المسكوفيت
قد تفكّك قبل ذلك تحت خط الصلابة، مع الكوارتز، إلى فلسبار بوتاسي
وسيليكات ألومنيوم، فغادر ماؤه مائعًا قبل بدء الانصهار؛ وقد تحوي \omterm{def:b2:magma-emplacement:magma}{الصهارة} ماءً
أكثر من \qty{6}{\percent}؛ وقد يبقى بعض الماء في البيوتيت.

\noindent\textbf{الجزء الثاني.}

\textbf{7.} $680 - 7.5\times(7.8 - 2) = \qty{636.5}{\celsius}$؛ وذروة الصخر تقع
تحته بمقدار \qty{43.5}{\celsius} فلا يستطيع أن ينصهر، حتى مع ماء حرّ.

\textbf{8.} $605 + (19.2 - 15.3)/(27 - 15.3)\times65 = \qty{626.7}{\celsius}$. عند
$t = 0$ تكون القاعدة، عند \qty{595}{\celsius}، تحته بمقدار \qty{32}{\celsius}؛ وعند
\qty{10}{Myr}، عند \qty{656}{\celsius}، تكون فوقه.

\textbf{9.} يحتاج الانصهار عند خط الصلابة الرطب إلى ماء حرّ، والصخر المتحوّل المسخَّن
لا يحوي منه إلا القليل جدًّا: فيُستهلك الماء الحرّ في \omterm{def:b2:magma-emplacement:magma}{الصهارة} الأولى، التي تتوقّف بعد ذلك عن النموّ
حتى تتفكّك الميكا.

\textbf{10.} بين 30 و\qty{40}{Myr}: $30 + (800 - 795)/(845 - 795)\times10 =
\qty{31.0}{Myr}$.

\textbf{11.} تنزل القاعدة تحت \qty{12.5}{kbar} بين 70 و\qty{80}{Myr}:
$70 + (13.7 - 12.5)/(13.7 - 12.4)\times10 = \qty{79.2}{Myr}$، عند
$914 - 14\times0.923 = \qty{901}{\celsius}$، داخل حصر الأمفيبول: فالأمفيبوليت
عند القاعدة سيكون في طور الانصهار.

\textbf{12.} $90 + (11.0 - 10)/(11.0 - 9.6)\times10 = \qty{97.1}{Myr}$، عند $867 -
53\times0.714 = \qty{829}{\celsius}$. وهذا فوق حصر المسكوفيت: فقد تفكّك المسكوفيت.
ويقع ضمن 800 إلى \qty{850}{\celsius} لتفكّك البيوتيت
(المقيس عند ضغط أدنى): فالبيوتيت، المعدن المائي الآخر في الشيست، عند
حدّه.

\noindent\textbf{الجزء الثالث.}

\textbf{13.} $n(\ce{Al2O3}) = 15.0/101.96 = \qty{0.14712}{mol}$، $n(\ce{CaO}) =
0.55/56.08 = 0.00981$، $n(\ce{Na2O}) = 3.4/61.98 = 0.05486$، $n(\ce{K2O}) = 4.9/94.20 =
0.05202$.

\textbf{14.} \mbox{\textdir TLT A/CNK $= 0.14712/0.11669 = 1.261$}.

\textbf{15.} \mbox{\textdir TLT A/NK $= 0.14712/0.10688 = 1.377$}؛ \mbox{\textdir TLT A/CNK $> 1$}: \omterm{def:b2:crustal-melting:asi}{فوق ألوميني}.

\textbf{16.} \omterm{def:b2:crustal-melting:asi}{فوق ألوميني}، ومن شيست: \omterm{def:b2:crustal-melting:types}{غرانيت من النمط S}. ويستقرّ ألومنيومه الفائض في
المسكوفيت والغارنت (أو الكورديريت).

\textbf{17.} $0.14712 - 0.11669 = \qty{0.03043}{mol}$ من \ce{Al2O3}؛ وفي صورة مسكوفيت،
$0.03043\times398.33 = \qty{12.1}{\gram}$ في \qty{100}{\gram} من الغرانيت.

\textbf{18.} الأوبسيديان \omterm{def:b2:crustal-melting:asi}{ميتاألوميني} (بالكاد)، دون الألومنيوم الفائض لصهارة
رسوبية متحوّلة: فمصدره مختلف على الأرجح. لكن عيّنة واحدة لا تثبت إلا القليل:
فلكل متتالية غرانيتويدية أعضاء فوق ألومينية باعتدال، والمتتالية تُقيَّم باتجاهها
كله.

\noindent\textbf{الجزء الرابع.}

\textbf{19.} $10^{12}\,\unit{\cubic\metre}\times2650 = \qty{2.65e15}{kg}$.

\textbf{20.} الغرانيت هو \omterm{def:b2:magma-emplacement:magma}{الصهارة} ناقص مائها: $2.65\times10^{15}/0.94 =
\qty{2.82e15}{kg}$ من \omterm{def:b2:magma-emplacement:magma}{الصهارة}.

\textbf{21.} لا يغادر بوصفه صهارة إلا $0.151\times0.70$ من كتلة الشيست:
$2.82\times10^{15}/(0.151\times0.70) = \qty{2.67e16}{kg}$.

\textbf{22.} $2.67\times10^{16}/2800 = \qty{9.5e12}{\cubic\metre}$، أي
\qty{9.5e3}{\cubic\kilo\metre}، نحو عشرة أضعاف \omterm{def:b2:magma-emplacement:bodies}{البلوتون}.

\textbf{23.} $9.5\times10^3/5 = \qty{1.9e3}{\square\kilo\metre}$، أي مربع طول ضلعه نحو
\qty{44}{km}.

\textbf{24.} مع $F = \qty{7.21}{\percent}$، يكون المصدر أكبر بمقدار $15.1/7.21 = 2.1$ مرة،
أي نحو \qty{2.0e4}{\cubic\kilo\metre}. وحجم قشرة المصدر المنصهرة لصنع
\omterm{def:b2:magma-emplacement:bodies}{البلوتون}، بنسبة \omterm{def:b2:magma-emplacement:magma}{الصهارة} في السؤال~5، هو
\textbf{\qty[text-series-to-math=true]{9.5e3}{\cubic\kilo\metre}}.
\end{solution}
